\subsection{Amplification on a prescribed interval}\label{cc:sec:amplification}

Of the two entries $\nu=E_1\cdot BE_2$ and $\eta=E_3\cdot BE_2$ of the
older gradient introduced in Lemma~\ref{cc:lem:frame}, the coupling
$\nu$ helps the preceding shear amplify a new wave, while $\eta$
changes its direction and eventually ends the growth. A small ratio
$\eta/\nu$ therefore allows a large gain before decay. In the model
case of Subsection~\ref{cc:sec:intro-waves} and the limit of an
infinitely strong preceding shear, the amplitude of a new wave solves
a linear system of two equations with explicit coefficients
(Lemma~\ref{cc:lem:strong-shear}). Up to errors measured by one number
$\delta$, the amplitude of a general new wave solves the same system
(Lemma~\ref{cc:uf:exact:lem:additive}). Its solution with growing
initial data peaks after a large gain (Lemma~\ref{cc:lem:model}), and
the wave follows this solution with relative errors that grow only
polynomially in the length of the interval
(Lemma~\ref{cc:lem:finite-comparison}). The couplings of the new wave
are estimated in Lemma~\ref{cc:lem:new-wave-couplings}, and
Proposition~\ref{cc:prop:finite-return} locates the ratio prescribed
for its successor after the amplitude peaks in the strong-shear limit.

\subsubsection{The new wave in the preceding frame}\label{cc:sec:comparison}
Let $B(t)$ be a smooth trace-free matrix, and let the preceding wave
$(k_p,\mathbf a_p)$, with $k_p\ne0$ and $\mathbf a_p\ne0$, satisfy
\eqref{cc:eq:wave-ode} with $D$ replaced by $B$. Let $(E_1,E_2,E_3)$ be
its frame, $B_{ij}=E_i\cdot BE_j$ the entries of $B$ in this frame, and
$\sigma_p=|k_p|\,|\mathbf a_p|$ its strength, as in
Lemma~\ref{cc:lem:frame}. Thus \eqref{cc:eq:frame} holds with the
functions $B_{12}$, $B_{32}$ and $\sigma_p$ in the places of
$\nu$, $\eta$ and $\sigma$. The new wave evolves over
\[
 M=B+\mathbf a_p\otimes k_p+E_c=B+\sigma_p E_2\otimes E_1+E_c,
\]
the sum of the older gradient, the gradient of the preceding wave at
the origin, and a continuous matrix $E_c(t)$ of errors. Its wave vector
$k$ solves the first equation of \eqref{cc:eq:wave-ode} with $D=M$ and
equals $\lambda E_3$, with $\lambda>0$, at the activation time $t_0$.
For this $k$, an \emph{amplitude vector} is a solution $\mathbf a$ of
the second equation in \eqref{cc:eq:wave-ode} with $D=M$ satisfying
$\mathbf a(t)\cdot k(t)=0$. Only these ordinary differential equations
enter this subsection: the matrices $B$ and $E_c$ are prescribed, and no
equation is assumed for them. Subsection~\ref{cc:uf:sec:later-wave}
applies the results below with $M$ the velocity gradient of the parent
at the origin, which is a fixed point of the flow.

Freeze the strength and the two couplings at the activation time,
\[
 \sigma=\sigma_p(t_0),\qquad \nu=B_{12}(t_0)>0,\qquad
 \eta=B_{32}(t_0)>0,
\]
and put
\[
 r=\sqrt{2\sigma\nu},\qquad S=\frac\sigma r=\sqrt{\frac{\sigma}{2\nu}},\qquad
 \varepsilon=\frac\eta\nu,\qquad \tau=r(t-t_0).
\]
Since $\nu/\sigma=1/(2S^2)$, the growth rate
$\sqrt{2\nu(\sigma-\nu)}$ found in Subsection~\ref{cc:sec:intro-waves}
equals $r(1-1/(2S^2))^{1/2}$; thus $r$ is this rate up to a factor
close to one when $S$ is large. The variable
$\tau$ measures time in units of the growth time $r^{-1}$, and $S$ is
the ratio of the preceding shear rate to the growth rate. We work on the
interval $0\le\tau\le T$, or $t_0\le t\le t_0+T/r$, with $T\ge2$.
In \eqref{cc:eq:wave-ode} and \eqref{cc:eq:frame} primes denote
derivatives in $t$. From now on in this subsection they denote
derivatives in $\tau$, so that $d/d\tau=r^{-1}d/dt$. All suprema are
over $0\le\tau\le T$, and $b_{ij}=B_{ij}/r$.

Define the scaled components $\zeta=(\zeta_1,\zeta_2,\zeta_3)$ and
$\alpha=(\alpha_1,\alpha_2,\alpha_3)$ of the wave vector and an amplitude
vector by
\begin{equation}\label{cc:eq:scaled-pair}
 k=\lambda\left(\zeta_1E_1+\frac{\zeta_2}{S}E_2+\zeta_3E_3\right),
 \qquad
 \mathbf a=\frac{\alpha_1}{S}E_1+\alpha_2E_2+\frac{\alpha_3}{S}E_3.
\end{equation}
Then $\zeta(0)=(0,0,1)$ and
\[
 |k|^2=\lambda^2m,\qquad
 m=\zeta_1^2+\zeta_2^2/S^2+\zeta_3^2.
\]
Transversality is $\zeta\cdot\alpha=0$, so where $\zeta_3\ne0$,
\[
 \alpha_3=-\frac{\zeta_1\alpha_1+\zeta_2\alpha_2}{\zeta_3}.
\]
The two components $\alpha_1,\alpha_2$ therefore determine the amplitude
vector. The factor $S^{-1}$ in the $E_2$ coordinate of $k$ makes the
contribution of $B$ to the coefficient of $\zeta_3$ in the equation for
$\zeta_2$ equal to $-B_{32}/\nu$, which equals $-\varepsilon$ at
$\tau=0$. In the growing
mode of Subsection~\ref{cc:sec:intro-waves}, the components of the
amplitude along $E_1$ and $E_2$ have opposite signs, and the first is
$S^{-1}(1-\nu/\sigma)^{-1/2}$ times the second in absolute value.
With the scaling in \eqref{cc:eq:scaled-pair}, this mode is
$\alpha_1=-\alpha_2$ to leading order. In
Subsection~\ref{cc:sec:scalar} the components $(\alpha_1,\alpha_2)$ are
divided further by the amplitude in the strong-shear limit, which gives
the normalized coordinates \eqref{cc:eq:normalized}.

In the model case of Subsection~\ref{cc:sec:intro-waves}, the entries
of $B$ in the frame are the constants $B_{12}=\nu$ and $B_{32}=\eta$,
all other entries vanish, and $E_c=0$. Put
\begin{equation}\label{cc:eq:reference-angle}
 \tan\varphi=\tfrac12\varepsilon\tau^2,\qquad
 0\le\varphi<\tfrac\pi2,\qquad
 \zeta_0=(\tan\varphi,-\varepsilon\tau,1).
\end{equation}
The next lemma derives the equations for the scaled components in the
model case and their limit as the preceding shear becomes infinitely
strong. In the limit the wave vector is $\lambda(\tan\varphi E_1+E_3)$,
so $\varphi$ is its angle from $E_3$ toward $E_1$, as in
Subsection~\ref{cc:sec:intro-waves}.

\begin{lemma}\label{cc:lem:strong-shear}
Suppose that on $0\le\tau\le T$ the entries of $B$ in the frame are
$B_{12}=\nu$, $B_{32}=\eta$ and $B_{ij}=0$ otherwise, and that $E_c=0$.
Then $\sigma_p=\sigma$, and the scaled components of the wave vector
satisfy
\[
 \zeta_1'=-\Bigl(1+\frac1{2S^2}\Bigr)\zeta_2,\qquad
 \zeta_2'=-\varepsilon\zeta_3,\qquad
 \zeta_3'=\frac{\varepsilon}{2S^2}\zeta_2 .
\]
Where $\zeta_3\ne0$, the components $\alpha_1,\alpha_2$ of every
amplitude vector satisfy
\[
 \begin{aligned}
 \alpha_1'&=-\alpha_2+\zeta_1\Pi,\\
 \alpha_2'&=-\Bigl(1-\frac1{2S^2}\Bigr)\alpha_1
   +\frac{\varepsilon}{2S^2}\,\frac{\zeta_1\alpha_1+\zeta_2\alpha_2}{\zeta_3}
   +\frac{\zeta_2}{S^2}\,\Pi,
 \end{aligned}
 \qquad
 \Pi=\frac{2\zeta_2\alpha_1+(\zeta_1+\varepsilon\zeta_3)\alpha_2}m .
\]
As $S\to\infty$ with $\varepsilon$ and $T$ fixed, $\zeta\to\zeta_0$
uniformly on $[0,T]$, and the coefficients of this system for
$(\alpha_1,\alpha_2)$ converge uniformly on $[0,T]$ to those of
\begin{equation}\label{cc:eq:model-system}
 \alpha_1'=-\cos^2\varphi(1-\varepsilon\tan\varphi)\alpha_2
                 -\varepsilon\tau\sin2\varphi\,\alpha_1,
 \qquad \alpha_2'=-\alpha_1 .
\end{equation}
\end{lemma}
\begin{proof}
In this proof write $k_i=k\cdot E_i$ and $a_i=\mathbf a\cdot E_i$, and
let a dot denote the derivative in $t$. The frame equations
\eqref{cc:eq:frame} read $\dot E_1=-\nu E_2$,
$\dot E_2=\nu E_1-\eta E_3$ and $\dot E_3=\eta E_2$, and their last
equation gives $d(\log\sigma_p)/dt=B_{33}=0$, so $\sigma_p=\sigma$. Thus
$M=\nu E_1\otimes E_2+\eta E_3\otimes E_2+\sigma E_2\otimes E_1$, and
$ME_1=\sigma E_2$, $ME_2=\nu E_1+\eta E_3$, $ME_3=0$. Since
$\dot k=-M^Tk$, we have $\dot k_i=-k\cdot ME_i+k\cdot\dot E_i$, hence
\[
 \dot k_1=-(\sigma+\nu)k_2,\qquad \dot k_2=-2\eta k_3,\qquad
 \dot k_3=\eta k_2 .
\]
Since $\dot{\mathbf a}=-M\mathbf a+p_{\rm w}k$ and
$M\mathbf a=\nu a_2E_1+\sigma a_1E_2+\eta a_2E_3$, also
$\dot a_i=-E_i\cdot M\mathbf a+p_{\rm w}k_i+\mathbf a\cdot\dot E_i$, hence
\[
 \dot a_1=-2\nu a_2+p_{\rm w}k_1,\qquad
 \dot a_2=-(\sigma-\nu)a_1-\eta a_3+p_{\rm w}k_2,\qquad
 p_{\rm w}=\frac{2(\nu k_1a_2+\sigma k_2a_1+\eta k_3a_2)}{|k|^2}.
\]
Substitute \eqref{cc:eq:scaled-pair}, divide by $r$ to pass to the time
$\tau$, and use
\[
 rS=\sigma,\qquad \frac{\nu}{\sigma}=\frac1{2S^2},\qquad
 \frac{\nu S}r=\frac12,\qquad \frac{\eta S}r=\frac\varepsilon2,\qquad
 \frac\eta\sigma=\frac{\varepsilon}{2S^2}.
\]
The equations for $k_i$ give those for $\zeta$, and
$\lambda Sp_{\rm w}/r=\Pi$. The equations for $a_1$ and $a_2$ then give
those for $\alpha_1=Sa_1$ and $\alpha_2=a_2$, after substituting
$\alpha_3=-(\zeta_1\alpha_1+\zeta_2\alpha_2)/\zeta_3$.

As $S\to\infty$, the coefficients of the system for $\zeta$ converge to
those of $\zeta_1'=-\zeta_2$, $\zeta_2'=-\varepsilon\zeta_3$,
$\zeta_3'=0$, whose solution with $\zeta(0)=(0,0,1)$ is $\zeta_0$. Both
systems are linear with bounded coefficients on $[0,T]$, so
$\zeta\to\zeta_0$ uniformly there by Gronwall's inequality. In
particular $\zeta_3\to1$ and $m\to1+\tan^2\varphi=\cos^{-2}\varphi$.
In the equation for $\alpha_1$ the coefficients of $\alpha_1$ and
$\alpha_2$ therefore converge to
\[
 \frac{2\zeta_1\zeta_2}m\to-2\varepsilon\tau\tan\varphi\cos^2\varphi
   =-\varepsilon\tau\sin2\varphi,\qquad
 -1+\frac{\zeta_1(\zeta_1+\varepsilon\zeta_3)}m
 \to-\cos^2\varphi(1-\varepsilon\tan\varphi),
\]
since $(\tan\varphi+\varepsilon)\tan\varphi\cos^2\varphi
=\sin^2\varphi+\varepsilon\sin\varphi\cos\varphi$. In the equation for
$\alpha_2$ every term except $-\alpha_1$ has a factor $S^{-2}$ and a
bounded coefficient, so the two coefficients converge to $-1$ and $0$.
\end{proof}

We call \eqref{cc:eq:model-system} the \emph{reference system}. The
vector $\zeta_0$ solves the limiting system $\zeta_1'=-\zeta_2$,
$\zeta_2'=-\varepsilon\zeta_3$, $\zeta_3'=0$ for the wave vector.

The errors of the comparison are measured by
\begin{equation}\label{cc:uf:exact:eq:additive-error}
 \delta=\left\|\frac{B_{12}}{\nu}-1\right\|_\infty
       +\left\|\frac{B_{32}}{\eta}-1\right\|_\infty
       +\frac{\|B\|_\infty}{r}+\frac{\|E_c\|_\infty}{\nu}.
\end{equation}
The first two terms are the relative variations of the two couplings on
the interval. The third compares the older gradient with the growth
rate; it is at least $\nu/r=1/(2S)$, as shown below, so it is small
only if the preceding shear is strong. The last compares the error
$E_c$ with the coupling $\nu$ rather than with $r$. The scaled
components multiply the $E_2$ component of $k$ and the $E_1$ and $E_3$
components of $\mathbf a$ by $S$, so in the time $\tau$ an entry of
$E_c$ enters their equations with a factor up to $S/r=1/(2\nu)$; see
the proof of Lemma~\ref{cc:uf:exact:lem:additive}.

Fix an absolute constant $C_0$, sufficiently large for the finitely
many estimates of this subsection, and assume
\begin{equation}\label{cc:eq:finite-smallness}
 C_0\varepsilon\le1,\qquad C_0T^{C_0}\delta\le\varepsilon .
\end{equation}
All other constants and exponents $C,c$ in this subsection are
absolute and do not depend on $C_0$; they may change from one estimate
to the next. The errors below have the form $CT^C\delta$, and by
\eqref{cc:eq:finite-smallness} such an error is at most $\varepsilon$
times any prescribed negative power of $T$ once $C_0$ is large.

Condition \eqref{cc:eq:finite-smallness} forces the preceding shear to
be strong. Indeed, $|B|\ge|E_1\cdot BE_2|$, and at $\tau=0$ this entry
is $B_{12}(t_0)=\nu$. Since $\nu/r=\sqrt{\nu/(2\sigma)}=1/(2S)$, the
third term of \eqref{cc:uf:exact:eq:additive-error} gives
$\delta\ge\nu/r=1/(2S)$. On the other hand,
\eqref{cc:eq:finite-smallness} gives
$\delta\le\varepsilon C_0^{-1}T^{-C_0}\le C_0^{-2}T^{-C_0}$. For
$C_0\ge2$ and $T\ge2$ it follows that
\begin{equation}\label{cc:eq:S-bounds}
 S\ge\frac1{2\delta}\ge\frac{C_0^2T^{C_0}}2\ge T^2,\qquad
 S^{-2}\le4\delta^2\le\delta,
\end{equation}
where the last inequality uses $\delta\le1/4$.

In the applications in Subsection~\ref{cc:uf:sec:later-wave} and
Section~\ref{cc:uf:sec:tower}, the parameter choice
\eqref{cc:uf:tower:eq:schedule} prescribes $\varepsilon$, and
$\varepsilon^{-1}$ and $T$ are bounded by powers of the logarithm of
the frequency of the preceding wave, while $\delta$ is at most a
negative power of that frequency times a power of its logarithm.
Item~2 of the proof of Theorem~\ref{cc:uf:tower:thm:finite} verifies
\eqref{cc:eq:finite-smallness} in this way.

For a general background and $E_c=0$, each coefficient in the proof of
the next lemma is a polynomial in $\zeta$, $1/\zeta_3$, $1/m$ and the quantities
\begin{equation}\label{cc:eq:coefficient-list}
 b_{11},\ b_{22},\ b_{33},\ b_{13},\ b_{31},\
 \frac{b_{21}}S,\ \frac{b_{12}}S,\ \frac{b_{23}}S,\ \frac{b_{32}}S,\
 2Sb_{12},\ 2Sb_{32},\ \frac{\sigma_p}\sigma,\ S^{-2}.
\end{equation}
The \emph{reference value} of such an expression is its value at
$\zeta=\zeta_0$, $2Sb_{12}=1$, $2Sb_{32}=\varepsilon$ and
$\sigma_p/\sigma=1$, with all other quantities in
\eqref{cc:eq:coefficient-list} equal to zero. By
Lemma~\ref{cc:lem:strong-shear}, it is the limit of the expression in
the model case as $S\to\infty$.

\begin{lemma}\label{cc:uf:exact:lem:additive}
Assume \eqref{cc:eq:finite-smallness}. The scaled components of the wave
vector satisfy
\begin{equation}\label{cc:eq:wave-vector-comparison}
 \sup|\zeta-\zeta_0|\le CT^C\delta\le\tfrac12,\qquad
 \sup\left|\frac{|k|^2\cos^2\varphi}{\lambda^2}-1\right|
       \le CT^C\delta\le\tfrac12.
\end{equation}
In particular
\begin{equation}\label{cc:eq:zeta-bounds}
 \tfrac12\le\zeta_3\le\tfrac32,\qquad |\zeta|\le3T^2 .
\end{equation}
The components $\alpha_1,\alpha_2$ of every amplitude vector satisfy
\[
 \frac{d}{d\tau}\binom{\alpha_1}{\alpha_2}
 =\left\{-\begin{pmatrix}
 \varepsilon\tau\sin2\varphi&\cos^2\varphi(1-\varepsilon\tan\varphi)\\
 1&0
 \end{pmatrix}+\mathcal E(\tau)\right\}\binom{\alpha_1}{\alpha_2},
 \qquad \sup|\mathcal E|\le CT^C\delta,
\]
where $\mathcal E$ is the same for all amplitude vectors.
\end{lemma}
\begin{proof}
The equations \eqref{cc:eq:wave-ode} and \eqref{cc:eq:frame} are in the
time $t$, and every coefficient in them is a rate of change in $t$.
In this proof primes denote derivatives in $\tau$. Since
$d/d\tau=r^{-1}d/dt$, each coefficient is divided by $r$,
which turns $B_{ij}$ into $b_{ij}$ and $\sigma_p$ into
$S\sigma_p/\sigma$. By \eqref{cc:uf:exact:eq:additive-error} and
$r/S=2\nu$,
\[
 |b_{ij}|\le\delta,\qquad
 |2Sb_{12}-1|=\left|\frac{B_{12}}{\nu}-1\right|\le\delta,\qquad
 |2Sb_{32}-\varepsilon|
 =\varepsilon\left|\frac{B_{32}}{\eta}-1\right|\le\delta .
\]
In the time $\tau$ the last equation of \eqref{cc:eq:frame} reads
$(\log\sigma_p)'=b_{33}$, and $\sigma_p=\sigma$ at $\tau=0$. Hence
\[
 \frac{\sigma_p(\tau)}{\sigma}
 =\exp\left(\int_0^\tau b_{33}(\tau')\,d\tau'\right),
\]
and $\sup|\sigma_p/\sigma-1|\le2T\delta$, because the integral is at
most $T\delta\le1$ in absolute value.

When $E_c=0$, the equation for $k$ and
\eqref{cc:eq:frame} give
\[
 \zeta_2'=-b_{22}\zeta_2-2Sb_{32}\zeta_3,\qquad
 \zeta_3'=-b_{33}\zeta_3-(b_{23}-b_{32})\zeta_2/S,
\]
\[
 \zeta_1'=-b_{11}\zeta_1
  -(\sigma_p/\sigma+b_{21}/S+b_{12}/S)\zeta_2
  -(b_{31}+b_{13})\zeta_3.
\]
In the model case these are the equations of
Lemma~\ref{cc:lem:strong-shear}. When $E_c\ne0$, the term $-E_c^Tk$ is added to the equation for $k$,
and the frame still evolves over $B$. In the time $\tau$, the
additional coefficient matrix for $\zeta$ is
\[
 -r^{-1}\operatorname{diag}(1,S,1)\,E_c^T
                      \operatorname{diag}(1,S^{-1},1),
\]
where $E_c^T$ is the transpose of $E_c$ written in the frame $(E_i)$
of the preceding wave, the matrix with entries $E_j\cdot E_cE_i$.
Since $S\ge1$, its norm is at
most $(S/r)|E_c|=|E_c|/(2\nu)\le\delta$.

Write $\zeta'=(A_0+A_1)\zeta$, where $A_0$ is the matrix of the
limiting system $\zeta_1'=-\zeta_2$, $\zeta_2'=-\varepsilon\zeta_3$,
$\zeta_3'=0$, and $A_1$ collects the remaining coefficients. The bounds
above give $\sup|A_1|\le CT\delta$. The vector $\zeta_0$ solves
$\zeta_0'=A_0\zeta_0$ with the same initial value as $\zeta$, and
$|\zeta_0|\le2T^2$. Since $A_0^3=0$, its propagator is
$\exp(uA_0)=\Id+uA_0+u^2A_0^2/2$, of norm at most $2T^2$ for
$0\le u\le T$. Therefore
\[
 \zeta(\tau)-\zeta_0(\tau)
 =\int_0^\tau e^{(\tau-\tau')A_0}A_1(\tau')\zeta(\tau')\,d\tau',
\]
so
\[
 \sup|\zeta-\zeta_0|
 \le CT^C\delta\,\bigl(1+\sup|\zeta-\zeta_0|\bigr).
\]
Condition \eqref{cc:eq:finite-smallness} gives the first bound of
\eqref{cc:eq:wave-vector-comparison}, hence \eqref{cc:eq:zeta-bounds}
and $m\ge1/4$. Since $m=\zeta_1^2+S^{-2}\zeta_2^2+\zeta_3^2$ and
$S^{-2}\le\delta$ by \eqref{cc:eq:S-bounds}, it also
gives $|m-\cos^{-2}\varphi|\le CT^C\delta$. Multiplication by
$\cos^2\varphi\le1$ proves the second bound.

To express the pressure coefficient in the amplitude equation for
$E_c=0$, put
\begin{align*}
 F_1={}&\zeta_1(b_{11}-b_{13}\zeta_1/\zeta_3)
       +\zeta_2(\sigma_p/\sigma+b_{21}/S-(b_{23}/S)\zeta_1/\zeta_3)\\
      &+\zeta_3(b_{31}-b_{33}\zeta_1/\zeta_3),\\
 F_2={}&\zeta_1(Sb_{12}-b_{13}\zeta_2/\zeta_3)
       +\zeta_2(b_{22}-(b_{23}/S)\zeta_2/\zeta_3)\\
      &+\zeta_3(Sb_{32}-b_{33}\zeta_2/\zeta_3).
\end{align*}
Then the pressure coefficient of \eqref{cc:eq:wave-ode} satisfies
\[
 \frac{\lambda S}{r}p_{\rm w}
   =\frac{2(F_1\alpha_1+F_2\alpha_2)}m.
\]
Substituting for $\alpha_3$ in the first two amplitude equations gives
a system for $(\alpha_1,\alpha_2)^T$ whose coefficient matrix $N$ has entries
\begin{align*}
 N_{11}&=-b_{11}+2b_{13}\zeta_1/\zeta_3+\boxed{2\zeta_1F_1/m},\\
 N_{12}&=\boxed{-2Sb_{12}}+2b_{13}\zeta_2/\zeta_3+\boxed{2\zeta_1F_2/m},\\
 N_{21}&=\boxed{-\sigma_p/\sigma}-(b_{21}-b_{12})/S
       +((b_{23}+b_{32})/S)\zeta_1/\zeta_3+2\zeta_2F_1/(S^2m),\\
 N_{22}&=-b_{22}+((b_{23}+b_{32})/S)\zeta_2/\zeta_3
                                      +2\zeta_2F_2/(S^2m).
\end{align*}
Only the boxed terms have nonzero reference values. Every other term
has one of $b_{11}$, $b_{22}$, $b_{13}$, the quotients $b_{ij}/S$ or
$S^{-2}$ as a factor, so its reference value is zero. In $F_1$ only
the term $\zeta_2\sigma_p/\sigma$, and in $F_2$ only $\zeta_1Sb_{12}$
and $\zeta_3Sb_{32}$, have nonzero reference values. Hence the
reference values of $F_1$, $F_2$ and $m$ are $-\varepsilon\tau$,
$(\tan\varphi+\varepsilon)/2$ and $\cos^{-2}\varphi$. Thus the
reference entries of $N$ are
\[
 -\varepsilon\tau\sin2\varphi,\qquad
 -1+\tan\varphi(\tan\varphi+\varepsilon)\cos^2\varphi
     =-\cos^2\varphi(1-\varepsilon\tan\varphi),\qquad -1,\qquad 0,
\]
which are the coefficients of \eqref{cc:eq:model-system}.

Each entry of $N$ is a polynomial of bounded degree in $\zeta_1$,
$\zeta_2$, $\zeta_3$, $1/\zeta_3$, $1/m$ and the quantities
\eqref{cc:eq:coefficient-list}. These are at most $CT^2$.
Each differs from its reference value by at most $CT^C\delta$: the
coordinates by \eqref{cc:eq:wave-vector-comparison}, the reciprocals
by a multiple of $|m-\cos^{-2}\varphi|+|\zeta_3-1|$, the
coefficients by the bounds at the beginning of the proof, and $S^{-2}$
by \eqref{cc:eq:S-bounds}. Hence each
entry of $N$ differs from its reference value by at most $CT^C\delta$.

For $E_c\ne0$, the amplitude equation contains the additional term
$-(\Id-2\frac{k}{|k|}\otimes\frac{k}{|k|})E_c\mathbf a$, of norm at
most $|E_c|\,|\mathbf a|$. By \eqref{cc:eq:scaled-pair},
\eqref{cc:eq:zeta-bounds} and \eqref{cc:eq:S-bounds},
$|\mathbf a|\le C|(\alpha_1,\alpha_2)|$, while
$\alpha_1=S\mathbf a\cdot E_1$ and $\alpha_2=\mathbf a\cdot E_2$.
In the time $\tau$ this adds a coefficient matrix of norm at most
$C(S/r)|E_c|\le C\delta$ to the reduced amplitude system.
\end{proof}

\subsubsection{The reference solution}\label{cc:sec:scalar}
At $\tau=0$, the reference system \eqref{cc:eq:model-system} is
$\alpha_1'=-\alpha_2$, $\alpha_2'=-\alpha_1$. Its growing solutions
have $\alpha_1=-\alpha_2$, the direction $-E_1/S+E_2$ of the amplitude
in \eqref{cc:eq:scaled-pair}. As $\varphi$ increases, the pressure
reduces the coefficient driving $\alpha_1$ from one to
$\cos^2\varphi(1-\varepsilon\tan\varphi)$ and introduces the term
$-\varepsilon\tau\sin2\varphi\,\alpha_1$. The driving coefficient
becomes negative when $\varepsilon\tan\varphi>1$.

The \emph{reference solution} is the solution of the reference system
with $(\alpha_1,\alpha_2)(0)=(-1,1)$, shown in
Figure~\ref{cc:fig:wave-turning}. A later wave is started with these
data in Subsection~\ref{cc:uf:sec:later-wave}. The second equation of
\eqref{cc:eq:model-system}, $\alpha_2'=-\alpha_1$, holds at every
$\tau$. Hence, wherever $\alpha_2>0$, the logarithmic derivative of
$\alpha_2$ is
\[
 Y=\frac{\alpha_2'}{\alpha_2}=-\frac{\alpha_1}{\alpha_2},
\]
and the first equation of \eqref{cc:eq:model-system} gives the Riccati
equation
\begin{equation}\label{cc:eq:model}
 Y'=\cos^2\varphi(1-\varepsilon\tan\varphi)
           -\varepsilon\tau\sin2\varphi\,Y-Y^2,\qquad Y(0)=1.
\end{equation}
Conversely, if $Y$ solves \eqref{cc:eq:model}, then
$(-Ya_{\rm ref},a_{\rm ref})$ is the reference solution, where
\[
 a_{\rm ref}(\tau)=\exp\left(\int_0^\tau Y(\tau')\,d\tau'\right).
\]
We call $a_{\rm ref}$ the \emph{reference amplitude}.

\begin{lemma}\label{cc:lem:model}
Let $0<\varepsilon\le1$. The solution $Y$ of \eqref{cc:eq:model}
exists for all $\tau\ge0$, satisfies $|Y|\le1$, and has the form
\begin{equation}\label{cc:eq:Y-ansatz}
 Y=(1+e)\cos\varphi-\tfrac14\varepsilon\tau\sin2\varphi,
 \qquad 0\le e\le\tfrac18\varepsilon.
\end{equation}
The reference amplitude satisfies
\begin{equation}\label{cc:eq:reference-amplitude}
 a_{\rm ref}(\tau)=\sqrt{\cos\varphi(\tau)}
       \exp\left(\int_0^\tau(1+e)\cos\varphi\,d\tau'\right).
\end{equation}
The function $Y$ has exactly one zero $\tau_m$, with
$2\le\varepsilon\tau_m\le3$. It is positive before this zero and
negative afterwards, so $a_{\rm ref}$ is largest at $\tau_m$, and
\[
 \log a_{\rm ref}(\tau_m)\ge\varepsilon^{-1/2}-1.
\]
\end{lemma}
\begin{proof}
Since $\tan\varphi=\varepsilon\tau^2/2$,
\[
 \varphi'=\varepsilon\tau\cos^2\varphi,\qquad
 (\log\cos\varphi)'=-\tfrac12\varepsilon\tau\sin2\varphi.
\]
We look for $Y$ in the form \eqref{cc:eq:Y-ansatz}, an ansatz with an
unknown function $e$. To motivate it, write $Y=\cos\varphi+w$, where
$\cos\varphi$ is the positive equilibrium of the equation
$Y'=\cos^2\varphi-Y^2$ with $\varphi$ frozen. Neglecting $w'$, $w^2$,
$\varepsilon\tau w\sin2\varphi$ and the term
$-\varepsilon\sin\varphi\cos\varphi$ of order $\varepsilon$,
\eqref{cc:eq:model} reduces to
$2w\cos\varphi=-\varepsilon\tau\sin\varphi\cos^2\varphi$, that is,
$w=-\tfrac14\varepsilon\tau\sin2\varphi=\tfrac12(\log\cos\varphi)'$.
The factor $1+e$ absorbs the neglected terms.

Integrating \eqref{cc:eq:Y-ansatz} and using the formula for
$(\log\cos\varphi)'$ gives \eqref{cc:eq:reference-amplitude}.
Substituting \eqref{cc:eq:Y-ansatz} into \eqref{cc:eq:model} gives
\[
 e'=\cos\varphi\bigl(-2e-e^2
                 +\tfrac12\varepsilon\sin\varphi\cos\varphi\bigr),
 \qquad e(0)=0.
\]
The right-hand side is nonnegative at $e=0$ and nonpositive at
$e=\varepsilon/8$, because $\sin\varphi\cos\varphi\le1/2$. Hence
$e$ remains in $[0,\varepsilon/8]$ and exists for all $\tau\ge0$.
This constructs the global solution $Y$.

Wherever $Y\ge0$, \eqref{cc:eq:model} gives $Y'\le1-Y^2$, because
$\cos^2\varphi(1-\varepsilon\tan\varphi)\le1$. Since $Y(0)=1$,
$Y(\tau)$ never exceeds one. Also
$\sin\varphi\cos\varphi\le\min\{1/2,\cot\varphi\}$, so for $\tau>0$,
\[
 \tfrac14\varepsilon\tau\sin2\varphi
 \le\min\{\varepsilon\tau/4,1/\tau\}\le\tfrac12.
\]
Equation \eqref{cc:eq:Y-ansatz} therefore gives $Y\ge-1/2$.

Write $Y=\cos\varphi(1+e-\tfrac12\varepsilon\tau\sin\varphi)$.
It is positive for $\varepsilon\tau\le2$, because $e\ge0$ and
$\sin\varphi<1$. For $\varepsilon\tau=3$, we have
$\tan\varphi=9/(2\varepsilon)\ge9/2$, hence $\sin\varphi>3/4$ and
$\tfrac12\varepsilon\tau\sin\varphi>9/8\ge1+e$. Thus $Y$ is negative
there. At any zero with $\varepsilon\tau\ge2$,
\[
 Y'=\cos^2\varphi(1-\varepsilon\tan\varphi)
    =\cos^2\varphi\bigl(1-(\varepsilon\tau)^2/2\bigr)<0.
\]
Consequently $Y$ has exactly one zero $\tau_m$, with
$2\le\varepsilon\tau_m\le3$, and has the stated signs on its two sides.

For $0\le\tau\le(2/\varepsilon)^{1/2}$, we have $\tan\varphi\le1$,
hence $\cos\varphi\ge2^{-1/2}$. This interval ends before
$\tau_m\ge2/\varepsilon$, where $a_{\rm ref}$ is largest. By
\eqref{cc:eq:reference-amplitude} and $e\ge0$,
\[
 \log a_{\rm ref}(\tau_m)
 \ge\log a_{\rm ref}\bigl((2/\varepsilon)^{1/2}\bigr)
 \ge2^{-1/2}(2/\varepsilon)^{1/2}-\tfrac14\log2
 \ge\varepsilon^{-1/2}-1.
 \qedhere
\]
\end{proof}

Relative to the reference solution, the reference system is triangular.
To compare a solution $(\alpha_1,\alpha_2)$ of
\eqref{cc:eq:model-system} with the reference solution
$(-Ya_{\rm ref},a_{\rm ref})$, we divide by the reference amplitude,
which removes the exponential growth common to the growing solutions,
and we separate the size of $\alpha_2$ from the deviation of the ratio
$\alpha_1/\alpha_2$ from its reference value $-Y$. We call
\begin{equation}\label{cc:eq:normalized}
 h=\frac{\alpha_2}{a_{\rm ref}},\qquad
 l=-\frac{\alpha_1+Y\alpha_2}{a_{\rm ref}\cos\varphi}
\end{equation}
the \emph{normalized coordinates} of $(\alpha_1,\alpha_2)$. The factor
$1/\cos\varphi$ in $l$ gives the equations for $h$ and $l$ below the
same factor $\cos\varphi$, which yields a propagator bound uniform in
$\varepsilon$ and in the length of the interval. By the second
equation of \eqref{cc:eq:model-system} and $a_{\rm ref}'=Ya_{\rm ref}$,
we have $h'=l\cos\varphi$. The trace of the reference coefficient matrix
is $-\varepsilon\tau\sin2\varphi=2(\log\cos\varphi)'$. By Liouville's
formula, the determinant
\[
 \det\begin{pmatrix}\alpha_1&-Ya_{\rm ref}\\
                      \alpha_2&a_{\rm ref}\end{pmatrix}
 =a_{\rm ref}(\alpha_1+Y\alpha_2)=-a_{\rm ref}^2l\cos\varphi
\]
is a constant multiple of $\cos^2\varphi$. Formula
\eqref{cc:eq:reference-amplitude} thus gives
\[
 l(\tau)=l(0)\exp\left(-2\int_0^\tau(1+e)\cos\varphi\,d\tau'\right).
\]
The reference system becomes
\begin{equation}\label{cc:eq:triangular-comparison}
 h'=l\cos\varphi,\qquad l'=-2(1+e)l\cos\varphi,
\end{equation}
and the reference solution is $(h,l)=(1,0)$. The second equation shows
that $|l|$ does not increase. Since $e\ge0$, for $\tau_1\le\tau_2$,
\[
 |h(\tau_2)-h(\tau_1)|\le\int_{\tau_1}^{\tau_2}|l|\cos\varphi
 \le\int_{\tau_1}^{\tau_2}(1+e)|l|\cos\varphi
 =\tfrac12\bigl(|l(\tau_1)|-|l(\tau_2)|\bigr).
\]
The triangle inequality and $|l(\tau_2)|\le|l(\tau_1)|$ give
\[
 |h(\tau_2)|+|l(\tau_2)|\le|h(\tau_1)|+|l(\tau_1)|.
\]
Since $|(h,l)|\le|h|+|l|\le\sqrt2\,|(h,l)|$, the propagator
$\Psi_0(\tau_2,\tau_1)$ of \eqref{cc:eq:triangular-comparison} has norm
at most $2$ for $\tau_1\le\tau_2$, uniformly in both times and in
$\varepsilon$.

\subsubsection{Comparison with the reference solution}
The next lemma shows that the wave follows the reference solution on
$0\le\tau\le T$. Dividing by $a_{\rm ref}$ removes the exponential
growth, and the bound for $\Psi_0$ gives errors polynomial in $T$.
A direct Gronwall estimate for $(\alpha_1,\alpha_2)$ would introduce
$\exp(CT)$, while \eqref{cc:eq:finite-smallness} controls polynomial
losses in $T$. The relative bounds give $|\mathbf a|$ up to a factor
$\exp(\pm\varepsilon)$, so the wave inherits the reference gain.
Recovering the normalized coordinates from an amplitude vector costs
a factor $S$, because its $E_1$ component is scaled by $S^{-1}$.

\begin{lemma}\label{cc:lem:finite-comparison}
Assume \eqref{cc:eq:finite-smallness}, and let $\mathbf a$ be the
amplitude vector with $\mathbf a(t_0)=-E_1(t_0)/S+E_2(t_0)$, that is,
with $(\alpha_1,\alpha_2)(0)=(-1,1)$ in \eqref{cc:eq:scaled-pair}, the
initial data of the reference solution.
Then $\alpha_2>0$ on $0\le\tau\le T$, and
\begin{equation}\label{cc:eq:amplitude-comparison}
 \left|\frac{\alpha_1}{\alpha_2}+Y\right|\le CT^C\delta,\qquad
 \alpha_2\le|\mathbf a|\le(1+CT^C\delta)\alpha_2.
\end{equation}
Moreover,
\begin{equation}\label{cc:eq:norm-rate-error}
 \left|\log\frac{|\mathbf a|}{a_{\rm ref}}\right|\le\varepsilon.
\end{equation}
For $0\le\tau_1\le\tau_2\le T$, let $\Phi(\tau_2,\tau_1)$ map
$\mathbf b(\tau_1)$ to $\mathbf b(\tau_2)$ for every amplitude vector
$\mathbf b$. It maps the plane perpendicular to $k(\tau_1)$ onto the
plane perpendicular to $k(\tau_2)$, and
\[
 \|\Phi(\tau_2,\tau_1)\|
 \le CS\,\frac{|\mathbf a(\tau_2)|}{|\mathbf a(\tau_1)|}.
\]
\end{lemma}
\begin{proof}
For an amplitude vector $\mathbf b$, let $\beta=(\beta_1,\beta_2,\beta_3)$
be its scaled components, defined as those of $\mathbf a$ in
\eqref{cc:eq:scaled-pair}:
\[
 \mathbf b=\frac{\beta_1}SE_1+\beta_2E_2+\frac{\beta_3}SE_3 .
\]
Its normalized coordinates are given by \eqref{cc:eq:normalized} with
$(\beta_1,\beta_2)$ in place of $(\alpha_1,\alpha_2)$, that is,
\[
 \binom hl=\frac1{a_{\rm ref}}R\binom{\beta_1}{\beta_2},\qquad
 R=\begin{pmatrix}0&1\\-(\cos\varphi)^{-1}&-Y(\cos\varphi)^{-1}\end{pmatrix},
 \qquad R^{-1}=\begin{pmatrix}-Y&-\cos\varphi\\1&0\end{pmatrix}.
\]
By Lemma~\ref{cc:uf:exact:lem:additive}, $(h,l)$ solves
\eqref{cc:eq:triangular-comparison} with the additional coefficient
matrix $R\mathcal ER^{-1}$. Since $|Y|\le1$ and
$(\cos\varphi)^{-1}\le T^2$, we have $|R|\le CT^2$, $|R^{-1}|\le C$
and $|R\mathcal ER^{-1}|\le CT^C\delta$.

Let $\Psi(\tau_2,\tau_1)$ be the propagator of this perturbed system.
Variation of constants gives
\[
 \Psi(\tau_2,\tau_1)=\Psi_0(\tau_2,\tau_1)
 +\int_{\tau_1}^{\tau_2}\Psi_0(\tau_2,\tau')
       (R\mathcal ER^{-1})(\tau')\Psi(\tau',\tau_1)\,d\tau'.
\]
Consequently $\sup|\Psi|\le2+2T\cdot CT^C\delta\sup|\Psi|$, and
\eqref{cc:eq:finite-smallness} gives $|\Psi(\tau_2,\tau_1)|\le4$.

For the prescribed initial amplitude $\mathbf a(t_0)$ the normalized
coordinates at $\tau=0$ are $(h,l)(0)=(1,0)$, a constant solution of the reference
system. Applying the preceding formula to this vector gives
\[
 \sup|(h-1,l)|\le2T\cdot CT^C\delta\cdot4\le CT^C\delta\le\tfrac12.
\]
In particular $h\ge1/2$, and the coordinate identities give
\[
 \alpha_2=a_{\rm ref}h>0,\qquad
 \frac{\alpha_1}{\alpha_2}+Y=-\frac{l\cos\varphi}{h},\qquad
 \left|\frac{\alpha_1}{\alpha_2}+Y\right|\le2|l|\le CT^C\delta.
\]
This is the first bound of \eqref{cc:eq:amplitude-comparison}, and
$|Y|\le1$ gives $|\alpha_1|\le2\alpha_2$. Since
$\alpha_3=-(\zeta_1\alpha_1+\zeta_2\alpha_2)/\zeta_3$, the bounds
\eqref{cc:eq:zeta-bounds} give $|\alpha_3|\le18T^2\alpha_2$. By
\eqref{cc:eq:scaled-pair},
\[
 |\mathbf a|^2=\alpha_2^2+S^{-2}(\alpha_1^2+\alpha_3^2).
\]
Since $S^{-2}\le\delta$ by \eqref{cc:eq:S-bounds}, this proves the
second bound of
\eqref{cc:eq:amplitude-comparison}. Thus
$\log(|\mathbf a|/a_{\rm ref})=\log h+\log(|\mathbf a|/\alpha_2)$
is at most $CT^C\delta$ in absolute value, and
\eqref{cc:eq:finite-smallness} gives \eqref{cc:eq:norm-rate-error}.

To bound $\Phi(\tau_2,\tau_1)$, note that for an arbitrary amplitude
vector $\mathbf b$,
$\beta_1=-a_{\rm ref}(Yh+l\cos\varphi)$, $\beta_2=a_{\rm ref}h$ and
$\beta_3=-(\zeta_1\beta_1+\zeta_2\beta_2)/\zeta_3$. Hence
$\mathbf b/a_{\rm ref}=hL_h+lL_l$, where
\[
 L_h=-\frac YSE_1+E_2+\frac{\zeta_1Y-\zeta_2}{S\zeta_3}E_3,\qquad
 L_l=\frac{\cos\varphi}{S}\left(-E_1+\frac{\zeta_1}{\zeta_3}E_3\right).
\]
By \eqref{cc:eq:zeta-bounds} and $S\ge T^2$, $|L_h|+|L_l|\le C$.
Conversely,
$h=(\mathbf b/a_{\rm ref})\cdot E_2$, and
$lL_l=\mathbf b/a_{\rm ref}-hL_h$. By
\eqref{cc:eq:wave-vector-comparison},
\[
 |L_l|=\frac{\cos\varphi}{S}
          \sqrt{1+\zeta_1^2/\zeta_3^2}\ge\frac cS,
 \qquad |h|+|l|\le CS|\mathbf b|/a_{\rm ref}.
\]
Together with $|\Psi|\le4$, these two coordinate bounds give
\[
 \|\Phi(\tau_2,\tau_1)\|
 \le CS\,\frac{a_{\rm ref}(\tau_2)}{a_{\rm ref}(\tau_1)}
 \le CS\,\frac{|\mathbf a(\tau_2)|}{|\mathbf a(\tau_1)|},
\]
where the last step uses \eqref{cc:eq:norm-rate-error}.
\end{proof}

\subsubsection{The couplings of the new wave}
Recall that the new wave evolves over the gradient
$M=B+\sigma_pE_2\otimes E_1+E_c$ of Subsection~\ref{cc:sec:comparison}.
The couplings $\nu_+,\eta_+$ of $M$ in the frame of the new wave drive
its successor. The first controls the sign of the pressure coefficient,
since
\[
 p_{\rm w}=\frac{2k\cdot M\mathbf a}{|k|^2}
          =2\nu_+\frac{|\mathbf a|}{|k|}.
\]
Subsection~\ref{cc:uf:sec:parent-step} uses this positivity throughout
the amplification interval.

\begin{lemma}\label{cc:lem:new-wave-couplings}
Assume \eqref{cc:eq:finite-smallness}, and let $\mathbf a$ be the
amplitude vector of Lemma~\ref{cc:lem:finite-comparison}. Put
$E_{1,+}=k/|k|$, $E_{2,+}=\mathbf a/|\mathbf a|$ and
$E_{3,+}=E_{1,+}\times E_{2,+}$, and define
\[
 \nu_+=E_{1,+}\cdot ME_{2,+},\qquad
 \eta_+=E_{3,+}\cdot ME_{2,+}.
\]
Then $\nu_+\ge\varepsilon\nu/2$ on $0\le\tau\le T$. Where
$\varepsilon\tan\varphi\ge1$, that is, for
$\sqrt2/\varepsilon\le\tau\le T$,
\begin{equation}\label{cc:eq:actual-ratio}
 \left|\frac{\nu_+}{\nu}-1\right|\le\varepsilon,\qquad
 \left|\frac{\eta_+/\nu_+}{\cos\varphi}-1\right|\le\frac12.
\end{equation}
\end{lemma}
\begin{proof}
Put
\[
 q=\frac{\alpha_1}{\alpha_2},\qquad
 z=\frac{\alpha_3}{\alpha_2}
   =-\frac{\zeta_1q+\zeta_2}{\zeta_3},\qquad n=\frac{|\mathbf a|}{\alpha_2}.
\]
Then
\[
 \frac{\mathbf a}{\alpha_2}=\frac qSE_1+E_2+\frac zSE_3,\qquad
 n^2=1+\frac{q^2+z^2}{S^2}.
\]
For $E_c=0$, the components of $M\mathbf a/(r\alpha_2)$ in the frame
$(E_1,E_2,E_3)$ of the preceding wave are $(f_1/S,f_2,f_3/S)$, where
\[
 f_1=b_{11}q+Sb_{12}+b_{13}z,\qquad
 f_2=(\sigma_p/\sigma+b_{21}/S)q+b_{22}+(b_{23}/S)z,\qquad
 f_3=b_{31}q+Sb_{32}+b_{33}z.
\]
In the same frame, the components of
$(k/\lambda)\times(\mathbf a/\alpha_2)$ are
\[
 \left(\frac{\zeta_2z}{S^2}-\zeta_3,
       \frac{\zeta_3q-\zeta_1z}{S},
       \zeta_1-\frac{\zeta_2q}{S^2}\right).
\]
Since $E_{1,+}=k/|k|$, $E_{2,+}=\mathbf a/|\mathbf a|$ and
$E_{3,+}=(k\times\mathbf a)/(|k||\mathbf a|)$,
\[
 \nu_+=\frac{k\cdot M\mathbf a}{|k||\mathbf a|},\qquad
 \eta_+=\frac{(k\times\mathbf a)\cdot M\mathbf a}{|k||\mathbf a|^2}.
\]
We compute these scalar products in the frame of the preceding wave,
pairing $k/\lambda$ and $(k/\lambda)\times(\mathbf a/\alpha_2)$ with
$M\mathbf a/(r\alpha_2)$. With $|k|=\lambda\sqrt m$,
$|\mathbf a|=n\alpha_2$ and $r/S=2\nu$ this gives
\begin{align}\label{cc:eq:coupling-rational}
 \nu_+/\nu&=\frac{2(\zeta_1f_1+\zeta_2f_2+\zeta_3f_3)}{\sqrt m\,n},\\
 \eta_+/\nu&=\frac{2}{\sqrt m\,n^2}
  \left\{\left(-\zeta_3+\frac{\zeta_2z}{S^2}\right)f_1
       +(\zeta_3q-\zeta_1z)f_2
       +\left(\zeta_1-\frac{\zeta_2q}{S^2}\right)f_3\right\}.
 \label{cc:eq:coupling-rational-end}
\end{align}
The matrix $E_c$ adds $E_{1,+}\cdot E_cE_{2,+}/\nu$ to $\nu_+/\nu$
and $E_{3,+}\cdot E_cE_{2,+}/\nu$ to $\eta_+/\nu$. Each is at most
$\delta$ in absolute value.

At the reference values,
\[
 \zeta=(\tan\varphi,-\varepsilon\tau,1),\qquad
 q=-Y,\qquad z=Y\tan\varphi+\varepsilon\tau,
\]
and $f_1=1/2$, $f_2=-Y$, $f_3=\varepsilon/2$, $n=1$ and
$m=\cos^{-2}\varphi$. The two expressions in
\eqref{cc:eq:coupling-rational}--\eqref{cc:eq:coupling-rational-end}
are therefore, respectively,
\[
 \sin\varphi+\varepsilon\cos\varphi+2\varepsilon\tau Y\cos\varphi,
 \qquad
 (\varepsilon\tan\varphi-1)\cos\varphi
    +\frac{2Y^2}{\cos\varphi}+2\varepsilon\tau Y\sin\varphi.
\]
In the second expression the terms $\varepsilon\tan\varphi\cos\varphi$,
$2Y^2/\cos\varphi$ and $2\varepsilon\tau Y\sin\varphi$ are large compared
with $\cos\varphi$ when $\varepsilon\tan\varphi$ is large, and
$\varepsilon\tan\varphi=\tfrac12(\varepsilon\tau)^2$ grows with $\tau$.
Their large parts cancel. To see this, write $Y=(1+e-u)\cos\varphi$
by \eqref{cc:eq:Y-ansatz}, with $u=\tfrac12\varepsilon\tau\sin\varphi$,
so that $2\varepsilon\tau\sin\varphi=4u$ and
$2u^2=\varepsilon\tan\varphi\sin^2\varphi$. Dividing the second
expression by $\cos\varphi$ gives
\begin{align*}
 &(\varepsilon\tan\varphi-1)+2(1+e-u)^2+4u(1+e-u)\\
 &\qquad=\varepsilon\tan\varphi-1
   +\bigl[2(1+e)^2-4u(1+e)+2u^2\bigr]+\bigl[4u(1+e)-4u^2\bigr]\\
 &\qquad=1+4e+2e^2+\bigl(\varepsilon\tan\varphi-2u^2\bigr).
\end{align*}
The terms $\pm4u(1+e)$ cancel, and the large parts of
$\varepsilon\tan\varphi$ and $2u^2$ cancel:
\[
 \varepsilon\tan\varphi-2u^2=\varepsilon\tan\varphi\,(1-\sin^2\varphi)
 =\varepsilon\sin\varphi\cos\varphi\le\tfrac12\varepsilon .
\]
Similarly,
$2\varepsilon\tau Y\cos\varphi
=2(1+e)\varepsilon\tau\cos^2\varphi
 -2\varepsilon\sin^2\varphi\cos\varphi$.
Thus the reference values of $\nu_+/\nu$ and $\eta_+/\nu$ are
\begin{equation}\label{cc:eq:reference-couplings}
 \begin{gathered}
 \sin\varphi+\varepsilon\cos\varphi\cos2\varphi
                   +2(1+e)\varepsilon\tau\cos^2\varphi,\\
 (1+4e+2e^2+\varepsilon\sin\varphi\cos\varphi)\cos\varphi.
 \end{gathered}
\end{equation}

Where $\varepsilon\tan\varphi\ge1$, we have $\cos\varphi\le\cot\varphi\le\varepsilon$, so
$1-\varepsilon^2\le\sin\varphi\le1$. Also $0\le e\le\varepsilon/8$ and
\[
 \varepsilon\tau\cos^2\varphi
 \le\varepsilon\tau\cot^2\varphi
 =\frac4{\varepsilon\tau^3}\le\sqrt2\,\varepsilon^2.
\]
The two reference values in \eqref{cc:eq:reference-couplings} are
$1+O(\varepsilon^2)$ and $(1+O(\varepsilon))\cos\varphi$.

Where $\varepsilon\tan\varphi\le1$, we have $\varepsilon\tau\le\sqrt2$, hence $Y>0$ by
Lemma~\ref{cc:lem:model}. The unsimplified reference value of
$\nu_+/\nu$ is at least $\sin\varphi+\varepsilon\cos\varphi
\ge\varepsilon$.

For the wave, in
\eqref{cc:eq:coupling-rational}--\eqref{cc:eq:coupling-rational-end},
the coefficients differ from their reference values by at most
$CT\delta$; $\zeta$ differs from $\zeta_0$ by at most $CT^C\delta$;
the terms with $S^{-2}$ in those formulas, in $m$ and in $n$, are at
most $CT^C\delta$; and $q$ differs from $-Y$ by at most $CT^C\delta$,
by \eqref{cc:eq:amplitude-comparison}. The formula for $z$ and the
bound $\zeta_3\ge1/2$ give the same bound for its change.
The numerators are polynomials of bounded degree in these quantities
and $1/\zeta_3$, which are at most $CT^2$, while the denominators
$\sqrt m\,n$ and $\sqrt m\,n^2$ are at least $1/2$.
Together with the contribution of $E_c$, the two ratios therefore
differ from their reference values by at most $CT^C\delta$, which is
at most $\varepsilon T^{-2}/8$ by \eqref{cc:eq:finite-smallness}.

In both ranges the reference value of $\nu_+/\nu$ is at least
$\varepsilon$; in the late range this uses $C_0\varepsilon\le1$.
Hence $\nu_+/\nu\ge\varepsilon/2$ on the whole interval. In the late
range, $\cos\varphi\ge T^{-2}$. Combining the perturbation bounds with
\eqref{cc:eq:reference-couplings} gives
\[
 \left|\frac{\nu_+}{\nu}-1\right|\le\varepsilon,\qquad
 \frac{\eta_+}{\nu_+}=(1+O(\varepsilon))\cos\varphi.
\]
The latter gives \eqref{cc:eq:actual-ratio} for $C_0$ sufficiently large.
\end{proof}

\subsubsection{The crossing}
After the peak of the reference solution at $\tau_m$, the new wave
becomes the shear for the next wave, whose coupling ratio
$\varepsilon_+$ is prescribed in advance. By
\eqref{cc:eq:actual-ratio}, the coupling ratio $\eta_+/\nu_+$ of the
new wave is comparable to $\cos\varphi$. It therefore reaches
$\varepsilon_+$ when $\cos\varphi$ is comparable to $\varepsilon_+$,
that is, when $\tan\varphi$ is comparable to $1/\varepsilon_+$, or
$\tau$ to $(\varepsilon\varepsilon_+)^{-1/2}$.
We call a time $\tau$ at which $\eta_+/\nu_+=\varepsilon_+$ a
\emph{crossing}: there the ratio reaches the threshold $\varepsilon_+$
prescribed for the next wave. The next wave is activated at a crossing. The reference
amplitude decreases after $\tau_m$, and the wave-vector length is
comparable to $\lambda/\cos\varphi$. The next proposition bounds their
product for the actual wave at a crossing, and shows that its amplitude
is small compared with its maximum during the first unit of normalized
time. The strength $\sigma_+=|\mathbf a||k|$ at the crossing and the
coupling $\nu_+$ supply the frozen parameters $\sigma,\nu$ for the next
wave.

\begin{proposition}\label{cc:prop:finite-return}
Assume \eqref{cc:eq:finite-smallness}, let $\mathbf a$ be the amplitude
vector of Lemma~\ref{cc:lem:finite-comparison}, with the couplings of
Lemma~\ref{cc:lem:new-wave-couplings}, and suppose that
$0<\varepsilon_+\le\varepsilon/32$ and
$T\ge2(\varepsilon\varepsilon_+)^{-1/2}$. Then the interval
$\tau_m\le\tau\le2(\varepsilon\varepsilon_+)^{-1/2}$, where
$\tau_m\ge2/\varepsilon$ is the zero of $Y$
(Lemma~\ref{cc:lem:model}), contains a crossing. At every crossing in
this interval the strength $\sigma_+=|\mathbf a(\tau)||k(\tau)|$ satisfies
\begin{equation}\label{cc:eq:general-return}
 C^{-1}(\varepsilon\varepsilon_+)^{-1/2}
 \le\frac{\sigma_+}{\lambda\max_{[0,T]}|\mathbf a|}
 \le C(\varepsilon\varepsilon_+)^{-1/2}.
\end{equation}
Moreover
\begin{equation}\label{cc:eq:gain}
 |\mathbf a(\tau)|\le\exp\bigl(3-\varepsilon^{-1/2}\bigr)
 \max_{[0,T]}|\mathbf a|\qquad(0\le\tau\le1).
\end{equation}
\end{proposition}
\begin{proof}
By Lemma~\ref{cc:lem:model}, $2/\varepsilon\le\tau_m\le3/\varepsilon$.
Since $\varepsilon_+\le\varepsilon/32$, we have
$(\varepsilon\varepsilon_+)^{-1/2}\ge4/\varepsilon$. Thus
$\tau_m<2(\varepsilon\varepsilon_+)^{-1/2}\le T$, and $[\tau_m,T]$ lies
in the range $\tau\ge\sqrt2/\varepsilon$ of \eqref{cc:eq:actual-ratio}.

For $\tau_m\le\tau\le T$, we have
$\tan\varphi=\varepsilon\tau^2/2\ge2/\varepsilon\ge1$, hence
\[
 \frac1{2\tan\varphi}\le\cos\varphi\le\frac1{\tan\varphi},\qquad
 \frac1{4\tan\varphi}\le\frac{\eta_+}{\nu_+}\le\frac3{2\tan\varphi},
\]
by \eqref{cc:eq:actual-ratio}. At $\tau_m$,
$\tan\varphi\le9/(2\varepsilon)$, so
$\eta_+/\nu_+\ge\varepsilon/18>\varepsilon_+$. At
$\tau=2(\varepsilon\varepsilon_+)^{-1/2}$, we have
$\tan\varphi=2/\varepsilon_+$, so
$\eta_+/\nu_+\le3\varepsilon_+/4<\varepsilon_+$.
Continuity therefore gives a crossing in this interval. At every
crossing in $[\tau_m,T]$ the displayed bounds give
$1/(4\varepsilon_+)\le\tan\varphi\le3/(2\varepsilon_+)$, so
$\cos\varphi$ is comparable to $\varepsilon_+$.

To estimate the strength at a crossing, we compare the amplitude there
with its maximum on $[0,T]$. By \eqref{cc:eq:norm-rate-error},
$|\mathbf a|$ lies between $\exp(-\varepsilon)a_{\rm ref}$ and
$\exp(\varepsilon)a_{\rm ref}$. The maximum of $a_{\rm ref}$ on
$[0,T]$ is attained at $\tau_m$, so for $\tau_m\le\tau\le T$,
$|\mathbf a(\tau)|/\max_{[0,T]}|\mathbf a|$ is bounded above and below
by absolute multiples of $a_{\rm ref}(\tau)/a_{\rm ref}(\tau_m)$.
Formula \eqref{cc:eq:reference-amplitude} gives
\[
 \frac{a_{\rm ref}(\tau)}{a_{\rm ref}(\tau_m)}
 =\left(\frac{\cos\varphi(\tau)}{\cos\varphi(\tau_m)}\right)^{1/2}
   \exp\left(\int_{\tau_m}^{\tau}(1+e)\cos\varphi\,d\tau'\right).
\]
The integral lies between zero and two, since
$\int_{\tau_m}^\infty\cos\varphi\,d\tau\le2/(\varepsilon\tau_m)\le1$.
Also $|k|$ is comparable to $\lambda/\cos\varphi$, by
\eqref{cc:eq:wave-vector-comparison}. Thus
$\sigma_+/(\lambda\max_{[0,T]}|\mathbf a|)$ is comparable to
$(\cos\varphi(\tau)\cos\varphi(\tau_m))^{-1/2}$.
Here $\cos\varphi(\tau_m)$ is comparable to $\varepsilon$, because
$2/\varepsilon\le\tan\varphi(\tau_m)\le9/(2\varepsilon)$, and at a
crossing $\cos\varphi$ is comparable to $\varepsilon_+$. This proves
\eqref{cc:eq:general-return}.

For $0\le\tau\le1$, the bound $Y\le1$ gives
$a_{\rm ref}(\tau)\le\exp(\tau)\le\exp(1)$, while
$\log a_{\rm ref}(\tau_m)\ge\varepsilon^{-1/2}-1$ by Lemma~\ref{cc:lem:model}. By \eqref{cc:eq:norm-rate-error} at such $\tau$ and at $\tau_m$,
$|\mathbf a(\tau)|\le\exp(1+\varepsilon)$ and
$|\mathbf a(\tau_m)|\ge\exp(\varepsilon^{-1/2}-1-\varepsilon)$.
Since $2\varepsilon\le1$, this gives \eqref{cc:eq:gain}.
\end{proof}
